\subsection{Formal Metric Definitions}

This section defines the system-level metrics used to evaluate procedural behavior. All metrics are computed at the episode level unless otherwise stated.

\subsubsection{Spatial Variance}

Spatial variance measures the distribution diversity of agents across the environment. Let $\{(x_i, y_i)\}_{i=1}^{n}$ denote the positions of $n$ agents. The mean position is:

\[
\bar{x} = \frac{1}{n}\sum_{i=1}^{n} x_i, \quad \bar{y} = \frac{1}{n}\sum_{i=1}^{n} y_i
\]

Spatial variance is defined as:

\[
V = \frac{1}{n}\sum_{i=1}^{n} [(x_i - \bar{x})^2 + (y_i - \bar{y})^2]
\]

Higher values of $V$ indicate more dispersed and structurally diverse agent configurations.

\subsubsection{System Stability}

Stability is measured as the temporal variance of a selected system observable $M$ (e.g., average agent density or interaction rate) across episodes. Given episode-level measurements $\{M_1, M_2, \dots, M_T\}$:

\[
S = \frac{1}{T} \sum_{t=1}^{T} (M_t - \bar{M})^2
\]

Lower values of $S$ indicate more stable long-term system behavior.

\subsubsection{Emergence}

Emergence is quantified using state-space coverage over the course of a run. Let $\mathcal{S}$ denote the set of discretized system states encountered across all episodes, and let $|\mathcal{S}|$ denote its cardinality. Emergence is measured as:

\[
E = \frac{|\mathcal{S}|}{|\mathcal{S}_{\max}|}
\]

where $|\mathcal{S}_{\max}|$ is the maximum possible state-space size under discretization. This metric captures the system’s ability to explore nontrivial configurations.

\subsubsection{Collapse Detection}

System collapse is detected when diversity and agent activity fall below predefined thresholds for consecutive episodes. Formally, collapse is detected at episode $t$ if:

\[
D_t < \epsilon_D \quad \text{and} \quad A_t < \epsilon_A
\]

for $n$ consecutive episodes, where $A_t$ denotes average agent movement magnitude and $\epsilon_D, \epsilon_A$ are empirically defined thresholds.

\subsubsection{Constraint-Level Metrics}

For the evolving-constraint condition, additional metrics are computed:

\begin{itemize}
    \item Constraint population size: $|C_t|$
    \item Average constraint lifetime
    \item Constraint turnover rate between episodes
    \item Average marginal fitness contribution $\Delta E$
\end{itemize}

These metrics characterize the dynamics of the constraint ecology itself.
